\section*{Chapter \ref{ch:m1:financing} --- Financing: Repo, Securities Lending and Prime Brokerage}

\begin{solution}{exo:m1:financing:1}
$L = 6$; wiped out by a fall of $1/6 = 16.7\%$; $r_E = 6 \times 6\% - 5
\times 4\% = 16\%$.
\end{solution}

\begin{solution}{exo:m1:financing:2}
Cash raised $0.97 \times 50 = \$48.5$ million; interest $48.5\times10^6
\times 0.042 \times 7/360 = \$39\,608$.
\end{solution}

\begin{solution}{exo:m1:financing:3}
The rebate is $4\% - 6\% = -2\%$: the short costs \$400\,000 a year to carry.
\end{solution}

\begin{solution}{exo:m1:financing:4}
Loan \$48. Call at $48/(1-0.30) = \$68.57$. At \$60 the equity is \$12
against a required $0.40\times60 = \$24$: \$12 a share.
\end{solution}

\begin{solution}{exo:m1:financing:5}
Assets \$2.85 billion, equity $0.5 - 0.15 = \$0.35$ billion: $L = 8.14$.
Back to 6: sell $5 \times 0.05 \times 3 = \$0.75$ billion. To 4: assets must
be \$1.4 billion, so sell \$1.45 billion, half the book.
\end{solution}

\begin{solution}{exo:m1:financing:6}
$1 - 2.5/4 = 37.5\%$ of the book. Return on equity: $0.12\% \times 40 =
4.8\%$ before, $0.12\% \times 25 = 3.0\%$ after --- and the forced sale will
have moved the very spreads the fund is long.
\end{solution}

\begin{solution}{exo:m1:financing:7}
Bisection gives about 0.129 per unit sold. In a linearised spiral each
round's fall is $\lambda(L-1)$ times the previous one, where $\lambda$ is the
impact: the series diverges when $\lambda \ge 1/(L-1) = 0.143$. The numerical
threshold is a little lower because a divergent series is not needed: it is
enough that the \emph{sum} of the falls reaches $1/L = 12.5\%$ within twelve
rounds starting from 5\%.
\end{solution}

\begin{solution}{exo:m1:financing:8}
(i) Static margin on a doubled price is 3.75\% of current value, not 7.5\%.
(ii) The relevant horizon is not one day but the time needed to liquidate a
concentrated position, several days during which the dealer's own selling
moves the price. (iii) The historical worst day was observed while the
client was buying; the move that matters is the one that occurs when it, and
every other dealer holding the same hedge, is selling: the past distribution
excludes exactly that event. (Also: five years of a rising stock say little
about its left tail.)
\end{solution}

\begin{solution}{pb:m1:financing:1}
\textbf{1.} $L = 28/4 = 7$; wiped out by a fall of 14.3\%.
\textbf{2.} $28/1.6 = \$17.5$ billion; margin \$1.75 billion.
\textbf{3.} $1.75/28 = 6.25\%$.
\textbf{4.} \$2.25 billion is free cash held elsewhere. No dealer can see the
total position or verify that the cash is unencumbered: each margins its
fifth as if it were the whole.
\textbf{5.} $28 \times 4.45\% = \$1.246$ billion.
\textbf{6.} $1.246 - 0.16 = \$1.086$ billion, 27.2\% of equity.
\textbf{7.} $1.086/28 = 3.88\%$.
\textbf{8.} Loss \$2.8 billion; equity \$1.2 billion; \omterm{def:m1:financing:leverage}{leverage} $25.2/1.2 =
21$.
\textbf{9.} \$2.8 billion is called. The family office pays \$2.25 billion:
\$450 million per dealer against \$560 million owed, leaving \$110 million
unpaid at each.
\textbf{10.} Loss $0.08 \times 25.2 = \$2.016$ billion; equity $-\$0.816$
billion: the family office is insolvent and \omterm{def:m1:financing:leverage}{leverage} has no meaning.
\textbf{11.} Hedge value $23.184/5 = \$4.64$ billion. Loss owed on the swaps
$560 + 403 = \$963$ million; held against it, $350 + 450 = \$800$ million.
\textbf{12.} Shortfall \$163 million plus $0.12 \times 4\,637 = \$556$
million of liquidation loss: \$720 million.
\textbf{13.} $163 + 0.03 \times 4\,637 = \$302$ million.
\textbf{14.} \$2.8 billion instead of \$1.75 billion.
\textbf{15.} $T = 50$ days; impact $0.7\times0.03\times\sqrt5 = 4.7\%$; risk
$0.03\sqrt{50/3} = 12.2\%$; discount $4.7 + 1.645\times12.2 = 24.8\%$, four
times the 6.25\% held.
\textbf{16.} Four other dealers are liquidating the same stocks at the same
time, so the volume available to each is a fraction of the total and the
impact adds up; and the formula's random walk has no drift, while here the
whole market knows that \$23 billion is for sale.
\textbf{17.} The client's \emph{total} position across dealers. Each would
have seen a holding of twenty-five days' volume, not five.
\textbf{18.} Its sales were made before the others' sales had moved the
price; the last to sell sold into the impact of everyone else. In a common
liquidation the order of exit is the whole result, which is why informal
agreements to sell in an orderly way do not hold.
\textbf{19.} \textbf{14.3\%.}
\textbf{20.} \omterm{def:m1:financing:leverage}{Leverage} measures how much the client loses per unit of price
move; margin measures how much of that loss the lender has been given in
advance.
\end{solution}

\begin{solution}{iq:m1:financing:1}
Obtain a \omterm{def:m1:financing:seclending}{locate} from the \omterm{def:m1:financing:pb}{prime broker}. Sell the shares in the market. At
settlement the \omterm{def:m1:financing:pb}{prime broker} borrows the shares from a lender and delivers
them to the buyer; the sale proceeds go to the lender as cash collateral,
marked to market daily. While short: receive interest on the collateral at
the \omterm{def:m1:financing:seclending}{rebate rate} (rate minus fee), pay the lender any dividend. To close: buy
the shares, return them to the lender, get the collateral back; profit is
the fall in price plus rebate less dividends.
\iqlookfor{that the seller never holds the proceeds, and the dividend.}
\end{solution}

\begin{solution}{iq:m1:financing:2}
A sale of a security with an agreement to repurchase it later at a higher
price: a collateralised cash loan. It is safe for the lender because it
holds a liquid security worth more than the loan (the \omterm{def:m1:financing:repo}{haircut}), marked to
market daily, which in most jurisdictions it can sell at once if the
borrower fails, without waiting for a bankruptcy court.
\iqlookfor{\omterm{def:m1:financing:repo}{haircut}, daily margining and the close-out right.}
\end{solution}

\begin{solution}{iq:m1:financing:3}
Sale $= (L-1)\,x\,A = 3 \times 5\% = 15\%$ of its assets. Check: assets
$100 \to 95$, equity $25 \to 20$, target assets 80, sell 15.
\iqlookfor{the $(L-1)x$ shortcut or a fast check with round numbers.}
\end{solution}

\begin{solution}{iq:m1:financing:4}
Financing spread on the long debit (around half a percent on the financed
part); loss of interest and credit spread on short proceeds; borrow fees,
negligible for large liquid names and many percent for crowded shorts, which
are often exactly the ones the model wants; dividends paid on shorts, if the
backtest used price returns; and withholding-tax effects on long dividends. For
a $3{:}2$ book on a unit of equity the first items alone are about 2\% of
equity a year; with hard-to-borrow shorts, several times that.
\iqlookfor{borrow fees correlated with the signal.}
\end{solution}

\begin{solution}{iq:m1:financing:5}
The client: no entry on the share register and, in some jurisdictions, no
disclosure of large holdings; no stamp or transaction tax where those apply;
negotiated margin; access to markets where it cannot hold shares directly.
The bank: a swap is a derivative on its balance sheet with different capital
and netting treatment, it earns the spread on the whole notional, and it
controls the hedge; against that, it carries the client's credit risk with
only the contractual margin, and a client's total position is harder to see.
\iqlookfor{both parties' motives, and the credit risk the dealer takes.}
\end{solution}

\begin{solution}{iq:m1:financing:6}
On liquidation cost, not on daily volatility: margin $\approx$ impact plus a
high quantile of the price move over the days needed to sell the position at
a realistic participation rate, with add-ons for concentration per name and
for the portfolio's lack of diversification. Make it dynamic, on current
value. Require disclosure of aggregate positions elsewhere as a condition of
the credit, and cap the limit if it is refused. Stress the whole portfolio
for a joint fall with correlated names. Make sure the contract allows margin
to be raised on short notice --- and use that right before the crisis, since
during it the client cannot pay.
\iqlookfor{days-to-liquidate as the horizon, and asking for the aggregate
position.}
\end{solution}
